\clearpage\pagestyle{otMain}\setlength{\otBlocEspacement}{2.5mm}\fontsize{11}{15}\selectfont\corrigesChapitre{15}{Droites, cercles et constructions}
\begin{blocLivre}{cor-15-1}\begin{corrige}{1}{Les bons noms}\demarche{Associer chaque symbole à ce qui est limité ou non.}
Le segment $\res{[AB]}$, la droite $\res{(CD)}$ et la demi-droite $\res{[EF)}$. Les crochets enferment les extrémités, les parenthèses laissent la droite filer.
\label{fig-cor-15-1}\figureGeo[corrige]{geo-1}
\end{corrige}\end{blocLivre}
\begin{blocLivre}{cor-15-2}\begin{corrige}{2}{Le milieu du segment}\demarche{Partager la longueur en deux parts égales.}
M est sur [AB], à égale distance des extrémités : AM $=7:2=\res{3{,}5}$ cm. Le corrigé montre M et le codage des deux moitiés égales.
\label{fig-cor-15-2}\figureGeo[corrige]{geo-2}
\end{corrige}\end{blocLivre}
\begin{blocLivre}{cor-15-3}\begin{corrige}{3}{Perpendiculaires ?}\demarche{Tester l’angle droit avec l’équerre, pas à l’œil.}
L’équerre s’ajuste exactement entre $(d_1)$ et $(d_2)$ : elles sont $\res{\text{perpendiculaires}}$. Les droites $(d_3)$ et $(d_4)$ se coupent sans angle droit : non perpendiculaires. Se couper ne suffit pas.
\label{fig-cor-15-3}\figureGeo[corrige]{geo-3}
\end{corrige}\end{blocLivre}
\begin{blocLivre}{cor-15-4}\begin{corrige}{4}{Parallèles ?}\demarche{Mesurer l’écart à deux endroits éloignés.}
L’écart entre $(a)$ et $(b)$ vaut 2 cm aux deux endroits mesurés : elles sont $\res{\text{parallèles}}$. L’écart entre $(a)$ et $(c)$ change : elles finiront par se couper.
\label{fig-cor-15-4}\figureGeo[corrige]{geo-4}
\end{corrige}\end{blocLivre}
\begin{blocLivre}{cor-15-5}\begin{corrige}{5}{Le vocabulaire du cercle}\demarche{Vérifier le passage par le centre.}
$\res{[OA]}$ est un rayon : il joint le centre au cercle. $\res{[BC]}$ est un diamètre : il passe par O. $[DE]$ est une corde : elle joint deux points du cercle sans passer par le centre.
\label{fig-cor-15-5}\figureGeo[corrige]{geo-5}
\end{corrige}\end{blocLivre}
\begin{blocLivre}{cor-15-6}\begin{corrige}{6}{Tracer une perpendiculaire}\demarche{Garder l’équerre appuyée sur la droite pendant le glissement.}
On pose un côté de l’angle droit de l’équerre sur (d), on glisse jusqu’à A, puis on trace. Le tracé, en vert au corrigé, forme un angle droit codé avec (d).
\label{fig-cor-15-6}\figureGeo[corrige]{geo-6}
\end{corrige}\end{blocLivre}
\begin{blocLivre}{cor-15-7}\begin{corrige}{7}{Tracer une parallèle}\demarche{Faire glisser l’équerre le long d’une règle fixe.}
On appuie l’équerre sur (d), la règle contre l’équerre, puis on fait glisser l’équerre jusqu’à B avant de tracer. La parallèle, en vert, garde un écart constant avec (d).
\label{fig-cor-15-7}\figureGeo[corrige]{geo-7}
\end{corrige}\end{blocLivre}
\begin{blocLivre}{cor-15-8}\begin{corrige}{8}{Rayon et diamètre}\demarche{Doubler ou partager en deux.}
Un diamètre vaut deux rayons : $4\times2=\res8$ cm. Dans l’autre sens : $7:2=\res{3{,}5}$ cm.
\label{fig-cor-15-8}\figureGeo[corrige]{geo-8}
\end{corrige}\end{blocLivre}
\begin{blocLivre}{cor-15-9}\begin{corrige}{9}{Lire les codes}\demarche{Faire parler chaque code, rien de plus.}
Les traits identiques affirment $\res{AB=AC}$ ; le petit carré affirme que l’angle en A est $\res{\text{droit}}$. Les codes énoncent les propriétés de la figure, même si le dessin est approximatif.
\label{fig-cor-15-9}\figureGeo[corrige]{geo-9}
\end{corrige}\end{blocLivre}
\begin{blocLivre}{cor-15-10}\begin{corrige}{10}{À 3 cm du puits}\demarche{Relier distance au centre et appartenance au cercle.}
Les points à exactement 3 cm forment le $\res{\text{cercle}}$ de centre O et de rayon 3 cm, tracé en vert. Les points à moins de 3 cm remplissent l’intérieur : c’est le $\res{\text{disque}}$ privé de son bord.
\label{fig-cor-15-10}\figureGeo[corrige]{geo-10}
\end{corrige}\end{blocLivre}
\begin{blocLivre}{cor-15-11}\begin{corrige}{11}{Un programme à exécuter}\demarche{Suivre les étapes dans l’ordre, instrument par instrument.}
Le rectangle se construit à l’équerre et à la règle. M est à $5:2=2{,}5$ cm de A ; le cercle de centre A passant par M a donc un rayon de $\res{2{,}5}$ cm. La construction complète figure en vert.
\label{fig-cor-15-11}\figureGeo[corrige]{geo-11}
\end{corrige}\end{blocLivre}
\begin{blocLivre}{cor-15-12}\begin{corrige}{12}{Deux cercles croisés}\demarche{Relire la définition du cercle pour conclure.}
Les cercles se coupent en deux points. Chacun est sur les deux cercles : il est à $\res{3\ \mathrm{cm}}$ de A et à $\res{3\ \mathrm{cm}}$ de B. Le compas garantit ces distances sans mesurer.
\label{fig-cor-15-12}\figureGeo[corrige]{geo-12}
\end{corrige}\end{blocLivre}
\begin{blocLivre}{cor-15-13}\begin{corrige}{13}{Reproduire sur quadrillage}\demarche{Compter les carreaux pour chaque sommet.}
On compte les carreaux : 5 vers la droite, 3 vers le haut pour les sommets, puis on joint. La diagonale relie deux sommets opposés. La reproduction figure en vert.
\label{fig-cor-15-13}\figureGeo[corrige]{geo-13}
\end{corrige}\end{blocLivre}
\begin{blocLivre}{cor-15-14}\begin{corrige}{14}{Perpendiculaires en chaîne}\demarche{Mesurer l’écart entre les deux droites extrêmes.}
$(d_1)$ et $(d_3)$ semblent ne jamais se couper : elles sont $\res{\text{parallèles}}$, l’écart mesuré reste constant. Deux droites perpendiculaires à une même droite sont parallèles entre elles.
\label{fig-cor-15-14}\figureGeo[corrige]{geo-14}
\end{corrige}\end{blocLivre}
\begin{blocLivre}{cor-15-15}\begin{corrige}{15}{Le plan du quartier}\demarche{Lire les codes du plan avant de répondre.}
Les rues des Lilas et des Érables gardent le même écart : $\res{\text{parallèles}}$. L’avenue du Parc coupe la rue des Lilas à angle droit codé : $\res{\text{perpendiculaires}}$. La fontaine est sur le cercle de centre F et sur celui de centre G de même rayon : on trace $\res{\text{deux arcs de cercle}}$ qui se croisent.
\label{fig-cor-15-15}\figureGeo[corrige]{geo-15}
\end{corrige}\end{blocLivre}
